\providecommand{\HMTParteIIIConstantesTransitionsRoot}{sections/transiciones_constantes}

\section{Transduction dimensionnelle des constantes engendrées en interne et comparaison métrologique ultérieure}

La classification généalogique du nombre clôt la publication mathématique : une section stable peut contracter son histoire jusqu'à une valeur sans perdre la mémoire qui détermine son identité. La mécanique dimensionnelle ajoute désormais une structure que le nombre, à lui seul, ne contient pas : représentation, bande, transducteur et persistance locale. La métrologie est la comparaison ultérieure entre cette réalisation et une carte dimensionnelle ; aucun étalon externe n'engendre rétrospectivement \(\alpha\), la décade \(10^{-34}\), Planck, le contre-angle, les constitutives du vide, la constante gravitationnelle ni le paramètre de Barbero--Immirzi.

\section{Transduction métrologique des états généalogiques : action, masse et observables dimensionnels}
\label{gmc:chap:programa}

Le pont serait superficiel s'il isolait les carrés de la thèse physique. HMT
ne présente pas APP, TRIT et TPK comme un jeu combinatoire auquel on
ajoute ensuite des constantes expérimentales. Son affirmation directrice est généalogique : les
constantes mathématiques et physiques sont des sorties de lecteurs construits sur la
même architecture, et la confrontation externe a lieu après la dérivation.
Le langage opératoriel doit rendre cette chaîne plus visible, non la réduire à un
ajustement numérique.

\subsection{La chaîne que le pont reçoit démontrée}

Dans les domaines et certificats de la théorie HMT--MD développée ici,
la construction reçoit comme
résultats démontrés :

\begin{enumerate}
 \item la génération monodromique de $\pi$, $e$, $\varphi$ et $\alpha$, avec
 la chaîne de neuf positions, ses étapes de filtration et les biographies de
 ses régions ;
 \item l'origine commune de la projection archimédienne et de la projection
 exceptionnelle, avec leurs codomaines différents ;
 \item l'action, la décade structurelle $10^{-34}$, la constante de Planck,
 le transducteur gravitationnel et le secteur Barbero--Immirzi ;
 \item les constitutives du vide et leur généalogie corrélée ;
 \item la Loi générale des masses sur le domaine complet
 $81/56/13/324$ et le secteur nul, sans les remplacer par une quelconque taxonomie expositive, suivie des réalisations physiques
 typées et des inventaires exhaustifs des masses et des largeurs, qui
 n'interviennent pas dans la construction du domaine ;
 \item la clôture électronique, les secteurs de mélange et de violation CP, et
 les transports CKM et PMNS avec leur statut mathématique spécifique dans
 chacun ;
 \item la sortie exceptionnelle, les codes et les symétries que l'œuvre
 déduit de la seconde projection.
\end{enumerate}

On n'écrit pas « sous des hypothèses déclarées » pour remplacer ces preuves. Les
hypothèses appartiennent à chaque théorème ; lorsqu'elles sont vérifiées, le
résultat est reçu comme démontré dans ce domaine. On n'exporte pas non plus vers un
domaine plus vaste une conclusion dont la réalisation physique requiert un morphisme
supplémentaire. Telle est la séparation probatoire du développement.

\subsection{Planck, action et la droite autoduale}

L'échelle absolue est séparée des rapports relatifs. Dans la notation de la
source des masses,
\[
 \Sigma_H=\varphi\alpha^{16},
 \qquad \operatorname{ord}_{10}(\Sigma_H)=-34,
 \tag{GMC-9.1}
\]
tandis que le facteur relatif d'action se construit dans une autre coordonnée. Dans
l'identité de périodicité 108--Compton--Planck, le point autodual satisfait, avec les
conventions de rayon déclarées,
\[
 R_{108}=\bar\lambda_C=r_g=\ell_P,
 \qquad C_{108}=2\pi\ell_P.
 \tag{GMC-9.2}
\]
La droite autoduale possède une conjugaison, une norme unitaire et une involution à point
fixe ; l'électron apparaît comme premier déplacement matériel, non comme origine
de l'échelle ni comme trou noir classique.

En langage opératoriel, $B_c$ produit l'intervalle libre
\textup{(GMC-6.15)} ; la rapidité $\eta$ de \textup{(GMC-2.6)} coordonne contraction et
ouverture hyperbolique ; l'application UCP publie un point de l'intervalle sans
épuiser la section qui l'a engendré. C'est un exemple complet de la
transduction recherchée : la nouvelle nomenclature ne change aucune identité
physique, mais sépare spectre, état, compression et généalogie.

\subsection{La loi des masses comme famille d'opérateurs positifs}

La loi HMT agit sur les fibres et les multisections au moyen d'opérateurs positifs.
Dans sa forme bidirectionnelle,
\[
 \widehat M_{M}^{\beta,r}
 =R_{\rm act}^{B_M^r/2}\,
  \widehat M_M^{(0)}\,
  R_{\rm act}^{B_M^r/2},
 \qquad r\in\{D,\Omega\}.
 \tag{GMC-9.3}
\]
La décade $10^{-34}$ fixe l'échelle absolue et s'annule dans les rapports ; les
lecteurs $K_D$ et $K_\Omega$ fixent le raffinement du chemin. Dans la fibre
électronique
\[
 K_D(\Gamma_e)=K_\Omega(\Gamma_e)=(80,54,6),
 \tag{GMC-9.4}
\]
et la réduction de rang un clôt la valeur électronique complète déterminée
par la clôture HMT. Les classes nulles, scalaires, opératorielles et non comparables
comme masse scalaire restent séparées : une matrice positive ne devient pas
un nombre expérimental avant la déclaration de la fonctionnelle de réalisation.

Le programme De las Cuevas--Netzer apporte ici une question nouvelle. Pour chaque
profondeur et niveau matriciel, quelles images UCP du système masse--chemin
conservent le spectre, lesquelles ne conservent que des moments et lesquelles se trouvent
dans l'enveloppe matricielle convexe de réalisations de rang un ? La
bigraduation $(m,s)$ transforme cette question en une famille de problèmes
finis cohérents.

\subsection{Matrice de transduction physico-mathématique}

\begin{longtable}{@{}L{.23\textwidth}L{.31\textwidth}L{.36\textwidth}@{}}
\toprule
Résultat HMT & Lectura operatorial & Question qui s'ouvre\\
\midrule
\endfirsthead
\toprule
Résultat HMT & Lectura operatorial & Question qui s'ouvre\\
\midrule
\endhead
monodromie $9\to1$ & retour de l'observable avec un état environnemental différent &
classer les dilatations qui préservent l'holonomie\\
$\pi,e,\varphi,\alpha$ & quatre lecteurs compatibles d'une généalogie &
existence et minimalité d'un système UCP commun\\
double projection & deux marginales d'un antécédent cylindrique &
antécédent non commutatif et séparation conjointe des fibres\\
$B_c$ et saut $4$ & intervalle matriciel libre $[5I,9I]$ &
extrémités et transport de la hiérarchie $4;2;6;54$\\
Planck et périodicité d'ordre 108 & point autodual et lecteurs conjugués &
réalisation comme symétrie d'un système conique avec mémoire\\
loi des masses & famille positive sur les fibres et les chemins &
rang matriciel, extrémités et réalisations par secteur\\
CKM/CP & transport de rang trois et phase non triviale &
canal covariant qui préserve spectre et orientation\\
symétrie exceptionnelle & second codomaine de la double projection &
entrelaceur entre atlas nonadique, code et représentation\\
\bottomrule
\end{longtable}

\subsection{Insertion universelle d'un signal APP dans un semi-groupe quantique de Markov}

Les constructions cycliques n'utilisent que le symbole latin. Il existe un cadre
général pour introduire tout signal hermitien des deux feuillets sans
rompre les sommes magiques.

\begin{proposition}[Double centrage d'un signal opératoriel]
\label{gmc:prop:doble-centrado}
Soit $G=(G_{ij})_{i,j=0}^{8}$ une grille de matrices autoadjointes dans
$M_3$. Définissons
\[
 H_{ij}=G_{ij}
 -\frac19\sum_{j'}G_{ij'}
 -\frac19\sum_{i'}G_{i'j}
 +\frac1{81}\sum_{i',j'}G_{i'j'}.
 \tag{GMC-9.5}
\]
Alors toutes les sommes de lignes et de colonnes de $H$ sont nulles. Si
$M=\max_{i,j}\|H_{ij}\|$ et si $\varepsilon\in\mathbb R$ satisfait
\[
 |\varepsilon|\le\frac1{9M}
 \tag{GMC-9.6}
\]
(avec la convention qu'il n'y a pas de restriction si $M=0$), la matrice
\[
 A_{ij}=\frac19I_3+\varepsilon H_{ij}
 \tag{GMC-9.7}
\]
est un \QMS de tailles $(9,3)$.
\end{proposition}

\begin{proof}
En sommant \textup{(GMC-9.5)} sur une ligne ou une colonne, les quatre termes s'annulent.
Par conséquent, \textup{(GMC-9.7)} a pour somme $I_3$ sur chaque ligne. Comme
$-MI_3\preceq H_{ij}\preceq MI_3$, la borne \textup{(GMC-9.6)} implique
$A_{ij}\succeq0$.
\end{proof}

Pour utiliser des sources de dimensions distinctes, il faut d'abord déclarer un lecteur
UCP vers $M_3$ : si $T_{ij}=T_{ij}^*$ appartient au système d'opérateurs
d'origine et si $\Theta$ est ce lecteur, on prend $G_{ij}=\Theta(T_{ij})$. On peut ainsi
insérer les images qutrit de signaux construits avec $P_\Sigma$,
$P_\Pi$, le bloc $B_c$ ou des lecteurs de chemin ; les opérateurs de tailles
$81\times81$ ou $2\times2$ ne sont pas introduits directement comme cellules de
$M_3$. Le double centrage
élimine toute composante purement « ligne seule + colonne seule » ; cela fournit
un falsificateur immédiat pour les codages qui présentent une apparence de couplage sans
avoir d'interaction bidimensionnelle.

\begin{controlbox}
La \Cref{gmc:prop:doble-centrado} démontre l'existence d'un \QMS marqué par un signal APP,
non sa fidélité ni son caractère non semi-classique. La faible amplitude peut le maintenir
dans un voisinage semi-classique. L'étape suivante doit résoudre le SDP
\textup{(GMC-8.6)} ou produire un témoin dual.
\end{controlbox}

\subsection{Factorisation du pont généalogique--matriciel au moyen d'immersions APP, d'un semi-groupe quantique de Markov, d'une frontière semi-classique, d'un canal continu, d'un système conique et d'une réalisation complètement positive}

Le pont laisse un programme ordonné par dépendances, non une liste ouverte
d'analogies :

\begin{enumerate}
 \item \textbf{Système APP filtré.} Construire des plongements complets
 d'ordre
 $\alpha_m:\mathcal S_m^{\rm APP}\hookrightarrow
 \mathcal S_{m+1}^{\rm APP}$ qui transportent les deux feuillets et la retenue.
 \item \textbf{QMS APP--TPK fidèle.} Choisir $G_{ij}$ dans
 \textup{(GMC-9.5)} à partir du couple somme--produit et prouver l'équation
 d'équivariance \textup{(GMC-7.9)}.
 \item \textbf{Frontière semi-classique.} Exploiter la symétrie nonadique pour
 réduire le SDP sur $S_9$ et décider si le candidat reste hors du système
 minimal.
 \item \textbf{Canal de continu.} Construire le tenseur local HMT et soumettre
 son canal de transfert à \textup{(GMC-7.2)} ; distinguer limite fine,
 grossière et mémoire des
 racines.
 \item \textbf{Système conique intrinsèque.} Définir le stem catégorique, les
 cônes, le dual et les transformations CP de TPK.
 \item \textbf{Réalisation métrologique conjointe.} Construire un système UCP
 commun aux lecteurs mathématiques, physiques et exceptionnels, et mesurer de quelle
 généalogie chaque sortie a besoin.
\end{enumerate}

\subsection{Hypothèses de naturalité, de cofinalité et de séparation requises par le théorème de reconstruction}

Une affirmation est établie lorsqu'elle inclut :

\begin{itemize}
 \item une définition complète du domaine et de tous les indices ;
 \item une preuve symbolique ou un certificat numérique avec tolérance déclarée ;
 \item une vérification de la positivité, de la normalisation et de la compatibilité ;
 \item des falsificateurs négatifs conservés ;
 \item une séparation entre coïncidence scalaire et équivalence de la tour
 matricielle ;
 \item une dérivation complète à partir des prémisses et des constructions qui la soutiennent.
\end{itemize}

Ce protocole n'atténue pas la thèse. Il permet de déclarer avec force ce qui
est déjà démontré et de localiser exactement l'extension capable de faire progresser
l'état de l'art.


L'électron constitue la première réalisation matérielle complète de ce passage. Son étape basale conserve sa fonction causale, mais la clôture bêta fixe la résidence directrice à partir de laquelle se développent les grandeurs, les champs et la Loi générale des masses.
